% ============================================================
% Section 6: Methodology
% ============================================================
\section{Methodology}

\subsection{End-to-End Workflow}
SafeHer operates on a continuous, event-driven loop following five sequential phases: (1)~Data Acquisition, (2)~Preprocessing, (3)~Risk Scoring, (4)~Decision \& Alerting, (5)~Feedback \& Logging.

\subsection{Phase 1: Contextual Data Collection}
The Flutter client collects sensor data at adaptive intervals based on motion state, as summarized in Table~\ref{tab:sensors}.

\begin{table}[htbp]
\caption{Sensor Data Collection Parameters}
\label{tab:sensors}
\centering
\scriptsize
\begin{tabularx}{\columnwidth}{|p{1.2cm}|p{1.8cm}|X|}
\hline
\textbf{Sensor} & \textbf{Sampling Rate} & \textbf{Purpose} \\
\hline
GPS & 30s (static) $\to$ 5s (moving) & Location, speed, route deviation \\
\hline
Accelerometer & 10Hz (burst on anomaly) & Fall detection, sudden stops \\
\hline
Gyroscope & 5Hz & Orientation changes, vehicle detection \\
\hline
System Clock & Event-triggered & Time-of-day, duration analysis \\
\hline
Network Info & On change & Connectivity risk factor \\
\hline
\end{tabularx}
\end{table}

Adaptive sampling reduces battery drain by 40--60\% compared to fixed-rate polling \cite{apple2024}.

\subsection{Phase 2: Feature Engineering \& Rule-Based Pre-Filtering}
Raw sensor data is transformed into contextual features including \texttt{time\_risk} (0--10 scale), \texttt{location\_type} (commercial/residential/isolated), \texttt{isolation\_index} (POI density within 200m), \texttt{velocity\_variance}, \texttt{route\_deviation}, \texttt{motion\_anomaly}, and \texttt{network\_risk}.

The rule-based risk aggregation applies deterministic thresholds:
\begin{lstlisting}[language=Python,caption={Rule-Based Risk Aggregation}]
Base_Risk = 0
IF time_risk >= 7:       Base_Risk += 15
IF location == "isolated": Base_Risk += 20
IF velocity_var > thresh AND motion_anomaly:
                         Base_Risk += 25
IF route_dev > 200m AND isolation > 0.8:
                         Base_Risk += 15
IF network_risk == 1:    Base_Risk += 10
Base_Risk = min(Base_Risk, 50)  # Cap at 50
\end{lstlisting}

\subsection{Phase 3: Hybrid Danger Score Computation}
The Verified Safety Matrix fuses rule-based and AI scores using weighted fusion:
\begin{equation}
S_{final} = \alpha \cdot S_{rule} + \beta \cdot S_{AI}
\label{eq:fusion}
\end{equation}
where $\alpha = 0.5$, $\beta = 0.5$ (calibrated via validation set), $S_{rule} \in [0, 50]$, $S_{AI} \in [0, 50]$, yielding $S_{final} \in [0, 100]$. Default alert threshold is 60, tunable per region/user via Firebase Remote Config.

\subsection{Phase 4: Autonomous Decision \& Alerting}
When $S_{final} \geq$ threshold and monitoring is not paused, the system constructs an alert payload containing user ID, timestamp, last known location, risk score, context summary, and a unique alert ID. It dispatches via Cloud Function, initiates a 15-minute live location stream, and logs to Firestore for audit.

\subsection{Phase 5: Feedback Loop \& Continuous Improvement}
Post-alert, trusted contacts can mark ``False Alarm'' or ``Confirmed Threat'' via the app UI. Anonymized feedback combined with context vectors is aggregated weekly to refine rule weights and AI prompt examples. A/B testing of thresholds and rule parameters is conducted across user cohorts to optimize precision/recall trade-offs.
